\chapter*{Заключение}
\addcontentsline{toc}{chapter}{Заключение}

В ходе работы работы был проведен анализ средств классификации и сегментации крюков в портале, 
объясненна их несостоятельность по отношению к крюкам и знаменам, и сделан вывод о необходимости 
нового решения для портала, в частности, в виде нейросетевой модели.

Были изучены правила постановки крюков и знамен в текстах древней руси, результаты показывают крайню 
сложноть в реализации любого логического алгоритма связи букв и крюков.

Анализ подходов к формированию выборки для нейросетевых моделей, сильно отразился на формировании итоговой
выборки, сделав ее объемной и сбалансированной

Разработанная модель для сегментации и классификации крюков отвечает поставленной задаче, по решению 
о запуске нейросетевой модели.

Разработанная модель интегрирована в портал Корпуса рукописного наследия древней руси.

Сформирована обширная выборка крюков и знамен для обучения нейросетевой модели.

%%% Local Variables:
%%% TeX-engine: xetex
%%% eval: (setq-local TeX-master (concat "../" (seq-find (-cut string-match ".*-3-pz\.tex$" <>) (directory-files ".."))))
%%% End:
